% Part: lambda-calculus
% Chapter: syntax
% Section: de-bruijn

\documentclass[../../../include/open-logic-section]{subfiles}

\begin{document}

\olfileid{lam}{syn}{deb}

\olsection{ڈی برائن دا اشاریہ}

$\alpha$-ہم ارزی بڑی سُبھاوک گل اے، کیوں جے
$\alpha$-ہم ارز ٹرماں دا «اکّو مطلب» اے۔ دراصل لامبڈا
ٹرماں دی ایہو جیہی نحو دِتّی جا سکدی اے جیہڑی اوہناں ٹرماں
وچ فرق نہیں کردی جنہاں دا $\alpha$-روپ اک دوجے وچ بدلیا جا
سکدا اے۔ ایس کم دی سبھ توں معروف نحو وچ متغیراں دی تھاں
اوہناں دا \emph{ڈی برائن اشاریہ} ورتیا جاندا اے۔

جدوں اسیں $\lambd[x][M]$ لکھدے آں، تاں صاف دسدے آں کہ
$x$ فنکشن دا پیرامیٹر اے، تاں جو $x$ نوں $M$ وچ اوس
پیرامیٹر دے حوالے لئی ورت سکیے۔ پر ڈی برائن اشاریے وچ
پیرامیٹراں دے ناں نہیں ہوندے۔ فنکشن دے اندر اوہناں دے حوالے
لئی اک عدد لکھیا جاندا اے، جیہڑا حوالے تے پیرامیٹر دے
وچکار تجریداں دیاں تہاں دی گنتی دسدا اے۔ مثال $\lambd[x][\lambd[y][y x]]$
ویکھو۔ باہرلی تجرید متغیر~$x$ نوں تے اندرلی تجرید متغیر~$y$
نوں پابند کردی اے۔ ذیلی ٹرم $y x$ اندرلی تجرید دے دائرۂ اثر
وچ اے۔ $y$ تے اوہدی تجرید~$\lambd[y]$ دے وچکار کوئی تجرید
نہیں، پر $x$ تے اوہدی تجرید~$\lambd[x]$ دے وچکار اک تجرید
اے۔ سو $0\, 1$ نوں $y x$ دی تھاں لکھدے آں، تے سارے ٹرم
لئی $\lambd[][\lambd[][01]]$ لکھدے آں۔

\begin{defn}
  ڈی برائن ٹرم استقرا نال ایویں متعین ہوندے نیں۔
  \begin{enumerate}
  \item $n$، جتھے $n$ کوئی وی قدرتی عدد اے۔
  \item $PQ$، جتھے $P$ تے $Q$ دوویں ڈی برائن ٹرم نیں۔
  \item $\lambd[][N]$، جتھے $N$ اک ڈی برائن ٹرم اے۔
  \end{enumerate}
\end{defn}

عام لامبڈا ٹرماں نوں ڈی برائن اشاریے والے ٹرماں وچ لکھن دا
باقاعدہ طریقہ ایویں اے۔
\begin{defn}
  \begin{align*}
    F_\Gamma(x) &= \Gamma(x) \\
    F_\Gamma(PQ) &= F_\Gamma(P)F_\Gamma(Q) \\
    F_\Gamma(\lambd[x][N]) &= \lambd[][F_{x,\Gamma}(N)]
  \end{align*}
  ایتھے $\Gamma$ متغیراں دی فہرست اے جیہدے اشاریے صفر توں
  شروع ہوندے نیں، تے $\Gamma(x)$ متغیر $x$ دی $\Gamma$ وچلی
  تھاں دسدا اے۔ مثلاً، جے $\Gamma$ دی فہرست $x,y,z$ اے، تاں
  $\Gamma(x)$ دی قدر $0$ تے $\Gamma(z)$ دی قدر $2$ اے۔

  $x,\Gamma$ اوہ فہرست اے جیہڑی $x$ نوں $\Gamma$ دے مُڈھ
  اُتے رکھ کے بندی اے۔ مثلاً، پچھلی مثال دی $\Gamma$ لئی
  $w,\Gamma$ دی فہرست $w,x,y,z$ اے۔
\end{defn}

ڈی برائن ٹرم توں عام لامبڈا ٹرم واپس حاصل کرن دا طریقہ ایویں اے۔

\begin{defn}
  \begin{align*}
    G_\Gamma(n) &= \Gamma[n] \\
    G_\Gamma(PQ) &= G_\Gamma(P) G_\Gamma(Q) \\
    G_\Gamma(\lambd[][N]) &= \lambd[x][G_{x,\Gamma}(N)]
  \end{align*}
  ایتھے $\Gamma$ فیر متغیراں دی فہرست اے جیہدے اشاریے صفر توں
  شروع ہوندے نیں، تے $\Gamma[n]$ تھاں~$n$ اُتے موجود متغیر دسدا اے۔
  مثلاً، جے $\Gamma$ دی فہرست $x,y,z$ اے، تاں $\Gamma[1]$ دا
  متغیر $y$ اے۔

  آخری مساوات وچ متغیر $x$ دی تھاں کوئی وی اوہ متغیر چُندے
  آں جیہڑا $\Gamma$ وچ نہیں۔
\end{defn}

ایتھے اسیں کجھ نتیجے ثبوت توں بغیر دیندے آں۔

\begin{prop}
  جے $M \aconvone M'$ تے $\Gamma$ کوئی وی اوہ فہرست اے جیہدے
  وچ $\FV{M}$ شامل اے، تاں $F_\Gamma(M) \eqs F_\Gamma(M')$۔
\end{prop}

\end{document}
